\documentclass[11pt]{article}

\usepackage[margin=1in]{geometry}
\usepackage{amsmath,amssymb,amsthm,bm}
\usepackage[T1]{fontenc}
\usepackage{lmodern}
\usepackage{microtype}
\usepackage{enumitem}
\usepackage{hyperref}

\title{HW1 -- Theory Solutions}
\author{\textit{Your Name}}
\date{\today}

\newcommand{\R}{\mathbb{R}}
\newcommand{\transpose}{^{\mathsf T}}
\newcommand{\loss}{\mathcal{L}}

\begin{document}
\maketitle

\section*{1. Normal Equations}

\begin{proof}
At a minimizer, the gradient of $\loss$ with respect to $\bm{\theta}$ is zero. Using $h_{\bm{\theta}}(\bm{x}_i)=\bm{\theta}\transpose\bm{x}_i$,
\begin{align*}
\nabla_{\bm{\theta}}\loss
&=\nabla_{\bm{\theta}}\left[\frac{1}{2M}\sum_{i=1}^{M}
	(\bm{\theta}\transpose\bm{x}_i-y_i)^2\right] \\
&=\frac{1}{M}\sum_{i=1}^{M}
	(\bm{\theta}\transpose\bm{x}_i-y_i)\bm{x}_i \\
&=\frac{1}{M}\bm{X}\transpose(\bm{X}\bm{\theta}-\bm{y}).
\end{align*}
Here, $\nabla_{\bm{\theta}}(\bm{\theta}\transpose\bm{x}_i-y_i)=\bm{x}_i$, since $\bm{x}_i$ and $y_i$ are fixed with respect to $\bm{\theta}$. Therefore,
\[
\bm{0}=\frac{1}{M}\bm{X}\transpose(\bm{X}\bm{\theta}-\bm{y})
\quad\Longrightarrow\quad
\bm{X}\transpose\bm{X}\bm{\theta}=\bm{X}\transpose\bm{y}.
\]
\end{proof}

\section*{2. Unique Solution}

\begin{proof}
We prove both directions.

\paragraph{If $\bm{X}$ has full column rank, the minimizer is unique.}
Full column rank means that $\bm{X}\bm{v}=\bm{0}$ only when $\bm{v}=\bm{0}$. Thus, for every nonzero $\bm{v}$,
\[
 \bm{v}\transpose\bm{X}\transpose\bm{X}\bm{v}
=(\bm{X}\bm{v})\transpose(\bm{X}\bm{v})
=\lVert\bm{X}\bm{v}\rVert^2>0,
\]
so $\bm{X}\transpose\bm{X}$ is invertible. Let
\[
\bm{\theta}^{*}=(\bm{X}\transpose\bm{X})^{-1}\bm{X}\transpose\bm{y}.
\]
Then $\bm{X}\transpose(\bm{X}\bm{\theta}^{*}-\bm{y})=\bm{0}$, so $\bm{\theta}^{*}$ satisfies the normal equations. For any $\bm{\theta}$, write $\bm{d}=\bm{\theta}-\bm{\theta}^{*}$. Using $\lVert\bm{a}+\bm{b}\rVert^2=\lVert\bm{a}\rVert^2+2\bm{a}\transpose\bm{b}+\lVert\bm{b}\rVert^2$,
\begin{align*}
\loss(\bm{\theta})-\loss(\bm{\theta}^{*})
&=\frac{1}{2M}\left(
 \lVert\bm{X}\bm{\theta}^{*}-\bm{y}+\bm{X}\bm{d}\rVert^2
 -\lVert\bm{X}\bm{\theta}^{*}-\bm{y}\rVert^2\right)\\
&=\frac{1}{2M}\left(
 \lVert\bm{X}\bm{\theta}^{*}-\bm{y}\rVert^2
 +2(\bm{X}\bm{\theta}^{*}-\bm{y})\transpose\bm{X}\bm{d}
 +\lVert\bm{X}\bm{d}\rVert^2
 -\lVert\bm{X}\bm{\theta}^{*}-\bm{y}\rVert^2\right)\\
&=\frac{1}{2M}\left(
 2(\bm{X}\bm{\theta}^{*}-\bm{y})\transpose\bm{X}\bm{d}
 +\lVert\bm{X}\bm{d}\rVert^2
\right)\\
&=\frac{1}{2M}\lVert\bm{X}\bm{d}\rVert^2\geq 0.
\end{align*}
The remaining cross term is zero because
\[
(\bm{X}\bm{\theta}^{*}-\bm{y})\transpose\bm{X}\bm{d}
=\bigl(\bm{X}\transpose(\bm{X}\bm{\theta}^{*}-\bm{y})\bigr)\transpose\bm{d}=0,
\]
Equality holds only if $\bm{X}\bm{d}=\bm{0}$, which by full column rank implies $\bm{d}=\bm{0}$. Therefore $\bm{\theta}^{*}$ is the unique minimizer.

\paragraph{If the minimizer is unique, $\bm{X}$ has full column rank.}
Suppose instead that $\bm{X}$ does not have full column rank. Then there is a nonzero vector $\bm{v}$ such that $\bm{X}\bm{v}=\bm{0}$. For a minimizer $\bm{\theta}^{*}$,
\[
 \loss(\bm{\theta}^{*}+\bm{v})
=\frac{1}{2M}\lVert\bm{X}(\bm{\theta}^{*}+\bm{v})-\bm{y}\rVert^2
=\frac{1}{2M}\lVert\bm{X}\bm{\theta}^{*}-\bm{y}\rVert^2
=\loss(\bm{\theta}^{*}).
\]
Since $\bm{v}\neq\bm{0}$, this gives a different minimizer, contradicting uniqueness. Hence $\bm{X}$ must have full column rank.
\end{proof}

\end{document}
